\documentclass[12pt,a4paper]{article}
\usepackage[UTF8]{ctex}
\usepackage{hyperref}
\usepackage{graphicx}
\usepackage{listings}
\usepackage{xcolor}
\usepackage{geometry}
\usepackage{longtable}
\usepackage{enumitem}
\geometry{margin=2.5cm}

\definecolor{codebg}{rgb}{0.95,0.95,0.95}
\lstset{
  backgroundcolor=\color{codebg},
  basicstyle=\ttfamily\small,
  breaklines=true,
  frame=single,
  columns=flexible,
  language=Python
}

\title{XENONnT Event Viewer \\ 软件设计、实现与优化完整文档}
\author{}
\date{2026年5月}

\begin{document}

\maketitle
\tableofcontents
\newpage

\section{项目概述}

\subsection{背景}

XENONnT 是位于意大利 Gran Sasso 地下实验室的暗物质探测实验。实验产生大量波形数据，研究人员需要交互式浏览工具来查看事件波形和 PMT 击中模式。EventViewer 是一个跨平台桌面应用，支持离线查看预提取的 NPZ 数据包。

\subsection{核心功能}

\begin{itemize}[itemsep=3pt]
\item \textbf{事件浏览}：左侧面板列出所有事件，支持 S1/S2 面积筛选、事件编号搜索
\item \textbf{Peak 列表}：可排序表格，显示 Type、Area (PE)、Width (ns)、Rise (ns)、Time (ns)
\item \textbf{主峰标注}：Main S1/Main S2 在列表中加粗显示
\item \textbf{三层波形显示}：Top PMT (蓝 \#2196F3) + Bottom PMT (绿 \#4CAF50) + Total (红 \#F44336)
\item \textbf{互动式 Peaks Zoom}：点击波形或表行进入单峰缩放模式
\item \textbf{图例交互切换}：点击图例文字切换波形层显示/隐藏
\item \textbf{滚轮缩放}：鼠标滚轮在波形面板上缩放
\item \textbf{页面缩放}：+/−按钮或 Ctrl+滚轮调整全局页面缩放百分比
\item \textbf{PMT ID 悬停}：鼠标悬停在 PMT 散点图上显示通道 ID
\item \textbf{PDF/PNG 导出}：Ctrl+S 直接保存到桌面
\item \textbf{Run Selector}：下拉菜单自动扫描切换内嵌 NPZ 数据文件
\item \textbf{多数据格式支持}：NPZ (peak\_basics) 模型波形，NPZ (real peaks) 真实波形
\end{itemize}

\subsection{技术栈}

\begin{itemize}[itemsep=2pt]
\item Python 3.9+
\item PySide6 (Qt GUI)
\item Matplotlib (科学绘图)
\item NumPy (数据处理)
\item Strax/Straxen (可选，用于直接数据访问)
\item PyInstaller (打包)
\end{itemize}

\section{架构设计}

\subsection{分层架构}

\begin{lstlisting}[language=text]
run_app.py                          # 入口点，命令行参数解析，自动加载 NPZ
+-- event_viewer_app/               # Qt GUI 层
|   +-- main_window.py              # MainWindow 主窗口
|   |   +-- PeakListWidget          # Peaks 列表组件
|   |   +-- Run Selector (QComboBox)# 运行选择器
|   +-- event_browser.py            # 事件列表，搜索，筛选
|   +-- event_canvas.py             # Matplotlib 画布 + 鼠标/键盘交互
|   +-- data_manager.py             # 数据加载与缓存管理
|   +-- qt_compat.py                # PySide6/PySide2 兼容层
+-- event_plotter/                  # 绘图引擎
|   +-- plotter.py                  # 核心绘图函数
|   |   +-- plot_event_full()       # 事件全览图
|   |   +-- plot_peak_zoom()        # 单峰缩放图
|   |   +-- plot_pmt_hit_pattern()  # PMT 击中模式
|   |   +-- plot_peaks()            # 全部 peaks 堆叠
|   |   +-- _draw_3layer_waveform() # 三层波形绘制
|   +-- io.py                       # 数据 I/O，PMT 几何加载
|   +-- style.py                    # 样式配置，颜色调色板
+-- scripts/                        # CLI 工具
|   +-- extract_event_bundle.py     # 从 strax 提取 NPZ
|   +-- plot_events.py              # 批量绘图
|   +-- build_mac.sh                # macOS PyInstaller 构建
+-- dali_probe/                     # DALI 数据探测
|   +-- extract_peaks_bundle.py     # DALI real peaks 提取
|   +-- report.md                   # DALI 探测报告
+-- debug_reports/                  # 调试报告
    +-- 00_SUMMARY.md               # 总结
    +-- 01_static_analysis.txt      # 静态分析
    +-- 02_data_pipeline_tests.txt  # 数据管线测试
    +-- 03_rendering_edge_cases.txt # 渲染边界测试
    +-- 04_ui_simulation_tests.txt  # UI 模拟测试
\end{lstlisting}

\subsection{数据流}

\begin{enumerate}[itemsep=3pt]
\item NPZ 文件 $\rightarrow$ \texttt{DataManager.open\_npz\_bundle()}
\item 加载 events[], peaks\_list[], eac\_list[], pmt\_positions, to\_pe
\item EventBrowser 填充事件列表，自动选中第一个事件
\item 用户点击事件 $\rightarrow$ \texttt{MainWindow.\_on\_event\_selected()}
\item $\rightarrow$ \texttt{plot\_event\_full()} $\rightarrow$ EventCanvas 渲染
\item Peak 点击 $\rightarrow$ EventCanvas 命中测试
\item $\rightarrow$ \texttt{peak\_clicked} Signal $\rightarrow$ \texttt{MainWindow.\_on\_peak\_clicked()}
\item $\rightarrow$ \texttt{plot\_peak\_zoom()} $\rightarrow$ 三面板视图
\end{enumerate}

\subsection{数据类型支持}

\begin{longtable}{|l|c|c|c|p{3cm}|}
\hline
\textbf{数据源} & \textbf{波形类型} & \textbf{PMT Pattern} & \textbf{data\_top} & \textbf{来源} \\
\hline
NPZ peak\_basics & 模型脉冲 (双指数) & EAC 按比例缩放 & 无 & Midway3 processed \\
NPZ real peaks & 真实 \texttt{data} 采样 & 真实 \texttt{area\_per\_channel} & 有 & DALI \\
Strax (cutax) & 真实 \texttt{data} & 真实 \texttt{area\_per\_channel} & 有 & 在线/离线 context \\
\hline
\caption{数据源对比}
\end{longtable}

\subsection{NPZ 数据格式}

NPZ bundle 包含以下键：
\begin{itemize}[itemsep=2pt]
\item \texttt{events}: 事件信息结构数组
\item \texttt{peaks\_list}: 每个事件的 peaks 列表（object dtype）
\item \texttt{eac\_list}: 每个事件的 event\_area\_per\_channel（可选）
\item \texttt{pmt\_x, pmt\_y, pmt\_array, pmt\_i}: PMT 位置（分列存储）
\item \texttt{to\_pe}: PMT 增益因子
\item \texttt{run\_id}: 运行编号
\item \texttt{waveform\_source}: ``real\_peaks'' 或 ``model''（新增字段）
\end{itemize}

\section{关键实现细节}

\subsection{画布管理：彻底解决残影}

\subsubsection{问题演化}

\begin{enumerate}[itemsep=3pt]
\item \textbf{初版 (fig.clear())}：重用 Figure，调用 \texttt{fig.clear()}。切换事件/peaks 时旧内容残留，出现"鬼影"重叠。
\item \textbf{fig.clf() 版}：改用 \texttt{fig.clf()} 完全清除。仍然有边缘彩色残影，尤其在滚动时。
\item \textbf{delaxes 版}：逐轴删除 \texttt{fig.delaxes(ax)}。残留依然存在。
\item \textbf{QScrollArea 版}：用 QScrollArea 包裹画布，\texttt{setWidgetResizable(False)}。边缘出现彩色像素点，页面显示"乱码"。
\item \textbf{最终方案}：每次 \texttt{clear()} 创建全新的 \texttt{Figure + FigureCanvasQTAgg}，旧画布 \texttt{deleteLater()}。彻底清除所有残留。
\end{enumerate}

\begin{lstlisting}[language=Python, caption={EventCanvas.clear() 最终实现}]
def clear(self):
    if self._hover_annot:
        try: self._hover_annot.remove()
        except Exception: pass
        self._hover_annot = None
    # Replace entire figure to eliminate ghosting
    old_canvas = self._canvas
    self._fig = Figure(figsize=(16, 10), facecolor="white", dpi=self._base_dpi)
    self._canvas = FigureCanvasQTAgg(self._fig)
    self._canvas.setSizePolicy(QSizePolicy.Expanding, QSizePolicy.Expanding)
    self._canvas.mpl_connect('button_press_event', self._on_click)
    self._canvas.mpl_connect('scroll_event', self._on_scroll)
    self._canvas.mpl_connect('motion_notify_event', self._on_hover)
    layout = self.layout()
    if layout:
        idx = layout.indexOf(old_canvas)
        if idx >= 0: layout.insertWidget(idx, self._canvas)
        else: layout.addWidget(self._canvas)
    old_canvas.setParent(None)
    old_canvas.deleteLater()
    if self._toolbar is not None:
        self._toolbar.canvas = self._canvas  # Fix save button
    self._update_canvas_size()
\end{lstlisting}

\subsection{交互系统}

\subsubsection{Peaks 点击命中测试}

点击波形区域后，通过邻近搜索找到最近的 peak：

\begin{lstlisting}[language=Python]
def _on_click(self, event):
    regions = getattr(event.inaxes, '_peak_regions', None)
    if not regions: return
    x = event.xdata
    MAX_CLICK_DIST = 0.0005  # 500 microseconds tolerance
    best, best_d = None, MAX_CLICK_DIST
    for r in regions:
        cx = r.get('center_x', (r['x_start'] + r['x_end']) / 2)
        d = abs(x - cx)
        if d < best_d: best_d = d; best = r
    if best is not None:
        self.peak_clicked.emit(best['index'])
\end{lstlisting}

每个 peak 在绘制时存储其时间区域到 \texttt{ax.\_peak\_regions} 列表。最初使用严格区间匹配（\texttt{x\_start <= x <= x\_end}），但窄峰（微秒级）在屏幕上不到1像素宽，无法有效点击。改为邻近搜索（500μs tolerance），同时保留区间匹配作为备选。

\subsubsection{图例交互切换}

三层波形的图例文字可点击切换每层的显示。实现方式：在 \texttt{\_on\_click} 中检测点击是否命中图例文字，若命中则切换对应层的可见性。

\begin{lstlisting}[language=Python]
def _check_legend_click(self, event):
    state = getattr(event.inaxes, '_legend_state', None)
    artists = getattr(event.inaxes, '_legend_artists', None)
    leg = event.inaxes.get_legend()
    for txt in leg.get_texts():
        if txt.contains(event)[0]:
            for key in state:
                if key in txt.get_text().lower():
                    state[key] = not state[key]
                    artists.get(key).set_visible(state[key])
                    return
\end{lstlisting}

\subsubsection{滚轮缩放}

鼠标滚轮在波形面板上缩放该面板的 x/y 轴范围。Page 缩放通过 Ctrl+滚轮或 +/- 按钮调整 Figure DPI。

\begin{lstlisting}[language=Python]
def _on_scroll(self, event):
    if event.inaxes is None: return
    s = 0.8 if event.button == 'up' else 1.25
    cx, cy = event.xdata, event.ydata
    xmin, xmax = ax.get_xlim()
    ymin, ymax = ax.get_ylim()
    ax.set_xlim(cx - (cx - xmin) * s, cx + (xmax - cx) * s)
    ax.set_ylim(cy - (cy - ymin) * s, cy + (ymax - cy) * s)
\end{lstlisting}

\subsubsection{PMT ID 悬停显示}

\texttt{motion\_notify\_event} 检测鼠标位置是否在 PMT 散点图的任一数据点上。若命中，显示黄色标注框显示 PMT 通道 ID。

\subsection{Peak 列表排序}

QTableWidget 默认按字符串排序，导致数值列（Area、Width 等）排序错误（如 ``9'' > ``100''）。

\textbf{解决方案}：自定义 \texttt{NumericItem} 类，重载 \texttt{\_\_lt\_\_} 方法实现数值比较。同时通过 \texttt{Qt.UserRole} 在每行保存原始 peak 索引，确保排序后点击正确的 peak。

\begin{lstlisting}[language=Python]
class NumericItem(QTableWidgetItem):
    def __lt__(self, other):
        try: return float(self.text()) < float(other.text())
        except: return super().__lt__(other)
\end{lstlisting}

\subsection{键盘快捷键}

方向键导航的难点：QListWidget 和 QTableWidget 会消费 Up/Down 按键。解决方案：

\begin{enumerate}[itemsep=2pt]
\item 使用 Left/Right 而非 Up/Down 切换事件
\item 使用 QShortcut 并指定 WindowShortcut 上下文：\\
  \texttt{QShortcut(QKeySequence("Left"), self).activated.connect(...)}
\item Escape 键同样使用 QShortcut
\end{enumerate}

\subsection{导出功能}

\subsubsection{演进}

\begin{enumerate}[itemsep=2pt]
\item \textbf{QFileDialog 版}：使用 \texttt{QFileDialog.getSaveFileName()} 弹出保存对话框。在 PyInstaller 打包的 .app 中对话框可能无法弹出（z-order 问题）。
\item \textbf{桌面直存版}：跳过对话框，直接保存到桌面，自动递增文件名避免覆盖。Ctrl+S 触发。
\item \textbf{工具栏保存按钮}：NavigationToolbar2QT 自带保存按钮，但画布替换后工具栏仍引用旧 canvas。修复：\texttt{self.\_toolbar.canvas = self.\_canvas}
\item \textbf{PDF 后端缺失}：PyInstaller 不包含 matplotlib.backends.backend\_pdf。修复：添加 \texttt{--hidden-import} 到构建参数。
\end{enumerate}

\subsection{PMT 击中模式渲染}

\subsubsection{圆圈半径修复}

硬编码 TPC 半径 47.9cm 远小于实际 PMT 阵列最大半径 65.6cm。初次实现时 252/494 个 PMT 在圆圈外。

\textbf{修复}：从 PMT 位置数据自动计算边界圆半径：
\begin{lstlisting}[language=Python]
pmt_r = np.sqrt(pmt_positions["x"]**2 + pmt_positions["y"]**2).max()
r_bound = pmt_r * 1.02
\end{lstlisting}

\subsubsection{颜色条 (Colorbar) 修复}

使用 \texttt{plt.colorbar()} 在画布替换后颜色条附加到错误的 Figure。全部替换为 \texttt{ax.figure.colorbar()}。

\subsubsection{PMT 点大小调整}

经历了多次迭代：50 → 25 → 12 → 25 → 60（最后一版加倍）。最终值：事件全览 120，Peaks Zoom 120。

\subsection{字体与样式优化}

经历了多轮调整：

\begin{longtable}{|l|c|c|c|}
\hline
\textbf{版本} & \textbf{基础字号} & \textbf{图例字号} & \textbf{坐标轴线宽} \\
\hline
初版 & 9pt & +0 (9pt) & 1.5 \\
第一次优化 & 11pt & +3 (14pt) & 2.0 \\
第二次优化 & 11pt & +6 (17pt) & 2.0 \\
最终版 & 11pt & +14 (25pt) & 2.0 \\
\hline
\caption{字体演化}
\end{longtable}

图例边框在最终版中完全移除（\texttt{legend.frameon = False}）。

标题从 \texttt{fontsize=\_rsize+1} 增加到 \texttt{+3}，suptitle 增加到 \texttt{+4}。

\subsection{波形线条粗细优化}

\begin{longtable}{|l|c|c|c|}
\hline
\textbf{组件} & \textbf{初版} & \textbf{第一次优化} & \textbf{最终版} \\
\hline
事件全览 Top/Bot & 0.25 & 0.75 & 0.75 \\
事件全览 Total & 0.35 & 1.05 & 1.05 \\
Peaks Zoom Top/Bot & 0.8 & 2.4 & 2.4 \\
Peaks Zoom Total & 1.5 & 4.5 & 4.5 \\
\hline
\caption{线宽演化}
\end{longtable}

所有线宽乘以 3 倍。数据模式（有波形采样）下的步进线宽从 [0.45, 0.45, 0.8] 增加到 [1.35, 1.35, 2.4]。

\subsection{波形颜色方案}

从柔和的 Nature 期刊配色（紫色/橙色/灰色）改为鲜艳的蓝绿红：

\begin{itemize}[itemsep=2pt]
\item Top PMT：\texttt{\#2196F3} (Material Blue) — 原为 Violet (\#9A4D8E)
\item Bottom PMT：\texttt{\#4CAF50} (Material Green) — 原为 Orange (\#E67E22)
\item Total：\texttt{\#F44336} (Material Red) — 原为 peak 类型色
\end{itemize}

\section{数据探索}

\subsection{数据源调查}

\subsubsection{DALI 探测}

DALI (\texttt{ssh dali}) 是 XENON 的专用数据节点。探测发现：
\begin{itemize}[itemsep=2pt]
\item \texttt{/dali/lgrandi/xenonnt/raw/}：仅有 raw\_records\_aqmon（监控数据），无 TPC 物理 raw\_records
\item \texttt{/dali/lgrandi/xenonnt/processed/}：包含带真实波形的 peaks（043864, 044116 等）
\item peaks dtype 包含：\texttt{data} (200样本, PE/sample)，\texttt{data\_top} (top阵列波形)，\texttt{area\_per\_channel} (494 PMT)
\end{itemize}

\subsubsection{Midway3 探测}

Midway3 (\texttt{ssh midway3}) 是主要计算集群：
\begin{itemize}[itemsep=2pt]
\item \texttt{/project/lgrandi/xenonnt/processed/}：2193 runs 的 peak\_basics
\item \texttt{/project2/lgrandi/xenonnt/processed/}：更多 runs
\item straxen.contexts.xenonnt()：查询了 73,925 个 run，0 个有 raw\_records 或 records 在本地存储
\item 023756 仅有 raw\_records\_aqmon，无 TPC raw\_records
\item 280 runs 有 peaklets（包含 data, area\_per\_channel）
\end{itemize}

\subsubsection{数据层级}

\begin{lstlisting}[language=text]
raw_records → records → peaklets → peaks → peak_basics
   (raw)      (ADC)     (per-PMT) (merged)   (metadata)
\end{lstlisting}

Midway3 主要存储：peak\_basics (元数据，轻量), event\_info, event\_area\_per\_channel。
DALI 额外存储：peaklets, peaks (带真实波形采样)。

\subsection{真实数据提取}

从 DALI 提取 real peaks 的流程：
\begin{enumerate}[itemsep=2pt]
\item SSH 到 DALI
\item source XENON 环境：\texttt{/cvmfs/xenon.opensciencegrid.org/releases/nT/el7.2025.07.2/setup.sh}
\item 运行 \texttt{dali\_probe/extract\_peaks\_bundle.py --run 043864 --n 200}
\item 生成 NPZ：9.1MB，200 events，3566 peaks
\item 从 Midway3 拉回本地（rsync，反复重试）
\end{enumerate}

\subsection{传输问题}

SCP 从 Midway3 传输大文件时常中断（9.1MB 文件经过多次尝试才完整拉回）。解决：
\begin{itemize}[itemsep=2pt]
\item 使用 rsync 替代 scp（支持断点续传）
\item 多次重试
\item 提取更小的样本（5-50 events）用于快速验证
\end{itemize}

\section{测试与验证}

\subsection{自动化测试套件}

独立调试报告位于 \texttt{debug\_reports/}：

\begin{longtable}{|l|c|c|c|}
\hline
\textbf{报告} & \textbf{测试数} & \textbf{结果} & \textbf{说明} \\
\hline
01\_static\_analysis & 11 files & 2 问题发现 & 导入、打印语句、硬编码路径 \\
02\_data\_pipeline & 9 项 & 全部通过 & 模式切换、缓存、索引 \\
03\_rendering\_edges & 8 项 & 全部通过 & 空 peaks、内存、DPI 极端值 \\
04\_ui\_simulation & 8 项 & 全部通过 & 排序、UserRole、主峰标记 \\
\hline
\caption{调试报告}
\end{longtable}

\subsection{端到端验证}

\begin{itemize}[itemsep=2pt]
\item 023756 NPZ (model): 50 events, 69 peaks/event — 全部加载
\item 043864 real peaks (5ev): 5 events, 真实波形 — 通过
\item 043864 real peaks (200ev): 200 events, 3566 peaks — 通过
\item 交叉切换：NPZ → real peaks → NPZ — 缓存正确清理
\item 20 次事件切换无渲染错误
\item 50 次 peak zoom 无 area 不一致
\item 100 次渲染无内存泄漏
\end{itemize}

\section{PyInstaller 打包}

\subsection{构建参数}

\begin{lstlisting}[language=bash]
pyinstaller \
  --onedir --windowed \
  --name XENONnT-EventViewer \
  --add-data "scripts/output/events_run_023756.npz:." \
  --hidden-import PySide6.QtCore \
  --hidden-import PySide6.QtGui \
  --hidden-import PySide6.QtWidgets \
  --hidden-import matplotlib.backends.backend_qtagg \
  --hidden-import matplotlib.backends.backend_pdf \
  --hidden-import matplotlib.backends.backend_agg \
  --collect-submodules numpy \
  --collect-data matplotlib \
  --exclude-module tkinter \
  --exclude-module PyQt5 \
  --exclude-module scipy \
  --exclude-module pandas \
  --distpath dist --workpath /tmp/pyinstaller_mac \
  run_app.py
\end{lstlisting}

\subsection{版本号}

构建脚本通过 PlistBuddy 注入版本号到 Info.plist：
\begin{lstlisting}[language=bash]
/usr/libexec/PlistBuddy -c "Set :CFBundleShortVersionString 2.0.0" \
  "dist/XENONnT-EventViewer.app/Contents/Info.plist"
\end{lstlisting}

\subsection{触发的 CI 构建}

GitHub Actions (\texttt{.github/workflows/build.yml}) 在推送 \texttt{v*} 标签时触发：
\begin{itemize}[itemsep=2pt]
\item build-linux: ubuntu-latest, PyInstaller → .tar.gz
\item build-windows: windows-latest, PyInstaller → .zip
\item build-macos-arm64: macos-latest, PyInstaller → .dmg
\item build-macos-x64: macos-13, PyInstaller → .dmg
\end{itemize}

\section{完整踩坑记录}

\begin{enumerate}[itemsep=5pt]

\item \textbf{QScrollArea 残影} \\
尝试方案：\texttt{setWidgetResizable(False/True)}、\texttt{setFixedSize}、\texttt{repaint()}、\texttt{draw()} vs \texttt{draw\_idle()}、scroll widget 嵌套。\\
最终方案：完全放弃 QScrollArea，每次 clear() 创建新 Figure+Canvas。

\item \textbf{工具栏保存按钮失效} \\
原因：画布替换后 NavigationToolbar2QT 仍引用旧 canvas。\\
修复：\texttt{self.\_toolbar.canvas = self.\_canvas}。

\item \textbf{PDF 导出缺少 backend} \\
原因：PyInstaller 不包含 matplotlib.backends.backend\_pdf。\\
修复：添加 \texttt{--hidden-import} 到构建脚本。

\item \textbf{Peak 列表排序错位} \\
原因：QTableWidget 字符串排序。\\
修复：NumericItem + Qt.UserRole 存储原始索引。

\item \textbf{macOS .app 打开终端} \\
原因：PyInstaller 单文件模式在 macOS 上打开终端。\\
修复：\texttt{--onedir --windowed}。

\item \textbf{strax 导入错误} \\
原因：strax 是可选依赖但被列在 hidden\_imports 中。\\
修复：从 hidden\_imports 移除 strax。

\item \textbf{PMT 圆圈不匹配} \\
原因：硬编码 47.9cm（TPC 半径）vs 实际 PMT 最大 65.6cm。\\
修复：从 PMT 位置自动计算。

\item \textbf{plt.colorbar 引用错误 Figure} \\
原因：\texttt{plt.colorbar()} 使用 pyplot 状态机，画布替换后引用错位。\\
修复：全部改为 \texttt{ax.figure.colorbar()}。

\item \textbf{内联注释破坏语法} \\
原因：全局查找替换颜色时，注释被插入函数调用中间。\\
样例：\texttt{color="\#2196F3"  \# bright blue, alpha\_fill=0.2} ← 注释吞掉了后续参数。\\
修复：逐行检查并移除函数参数中的内联注释。

\item \textbf{DALI/Midway3 SSH 传输不稳定} \\
原因：SCP 对大文件 (9MB+) 经常断连。\\
修复：rsync + 断点续传 + 多次重试。

\item \textbf{三层波形 PMT Pattern 全为零} \\
原因：NPZ 模式下 EAC 加载被 strax 模式守卫阻塞。\\
修复：移除 \texttt{if self.\_dm.mode == "strax"} 守卫，NPZ 模式也加载 EAC。

\item \textbf{事件波形空白（Build 不含 Codex 修改）} \\
原因：App 在 Codex 修改源码后构建，但源码修改未写入构建缓存。\\
修复：清除 PyInstaller 缓存，确认源文件时间戳，重新构建。

\item \textbf{滚轮缩放导致"乱码"} \\
原因：Trackpad 惯性滚动产生大量 scroll 事件，\texttt{draw\_idle()} 渲染跟不上。\\
修复：改为同步 \texttt{draw()} 渲染。

\item \textbf{切换事件/peaks 时旧波形残留} \\
原因：matplotlib 的 \texttt{fig.clear()} 和 \texttt{fig.clf()} 不会立即更新 Qt widget。\\
修复：最终采用全新 Figure+Canvas 方案。

\item \textbf{左侧面板不可调整大小} \\
原因：QSplitter stretchFactor(0,0) + setMaximumWidth(320) 锁定左侧面板。\\
修复：stretchFactor(0,1)，maximumWidth 改为 500。

\end{enumerate}

\section{目前已知限制}

\begin{itemize}[itemsep=3pt]
\item 023756 无 TPC raw\_records，波形为模型生成
\item real peaks 数据需要从 DALI 手动提取并传输
\item Strax 模式在主线程同步加载，大 run 可能卡住界面
\item Windows/Linux 版本需要 CI 或本地构建，无法在 macOS 上交叉编译
\item .app 使用 ad-hoc 签名，分发时可能触发 Gatekeeper 警告
\item Run Selector 扫描路径在打包后可能找不到外部 NPZ 文件
\end{itemize}

\section{参考文献}

\begin{itemize}[itemsep=2pt]
\item XENONnT Straxen 文档：\url{https://straxen.readthedocs.io}
\item PySide6 文档：\url{https://doc.qt.io/qtforpython-6/}
\item Matplotlib 文档：\url{https://matplotlib.org/stable/}
\item PyInstaller 手册：\url{https://pyinstaller.org/en/stable/}
\end{itemize}

\end{document}
